%% uncomment to list all files in log
%\listfiles

\documentclass[12pt]{report}


\usepackage{fontspec}

%\setmainfont[Scale=MatchLowercase]{Lucida Bright}
%\setmonofont{FreeMono}
%\setmonofont{Source Code Pro}
\setmonofont[Scale=MatchLowercase]{Ubuntu Mono}

\usepackage[headings]{fullpage}

% national use characters 
%\usepackage{inputenc}

% ams mathematical symbols
\usepackage{amsmath,amssymb}

% added to support pandoc highlighting
\usepackage{microtype}

\usepackage{makeidx}

% add index and bibliographies to table of contents
\usepackage[nottoc]{tocbibind}

% postscript courier and times in place of cm fonts
%\usepackage{courier}
%\usepackage{times}

% extended coloring
\usepackage{color}
\usepackage[table,dvipsnames]{xcolor}
\usepackage{colortbl}

% advanced date formating
\usepackage{datetime}

%support pandoc code highlighting
\usepackage{fancyvrb}
\DefineShortVerb[commandchars=\\\{\}]{\|}
\DefineVerbatimEnvironment{Highlighting}{Verbatim}{commandchars=\\\{\}}
% Add ',fontsize=\small' for more characters per line

%tango style colors
% \usepackage{framed}
% \definecolor{shadecolor}{RGB}{255,255,255}
% \newenvironment{Shaded}{\begin{snugshade}}{\end{snugshade}}
% \newcommand{\KeywordTok}[1]{\textcolor[rgb]{0.13,0.29,0.53}{\textbf{{#1}}}}
% \newcommand{\DataTypeTok}[1]{\textcolor[rgb]{0.13,0.29,0.53}{{#1}}}
% \newcommand{\DecValTok}[1]{\textcolor[rgb]{0.00,0.00,0.81}{{#1}}}
% \newcommand{\BaseNTok}[1]{\textcolor[rgb]{0.00,0.00,0.81}{{#1}}}
% \newcommand{\FloatTok}[1]{\textcolor[rgb]{0.00,0.00,0.81}{{#1}}}
% \newcommand{\CharTok}[1]{\textcolor[rgb]{0.31,0.60,0.02}{{#1}}}
% \newcommand{\StringTok}[1]{\textcolor[rgb]{0.31,0.60,0.02}{{#1}}}
% \newcommand{\CommentTok}[1]{\textcolor[rgb]{0.56,0.35,0.01}{\textit{{#1}}}}
% \newcommand{\OtherTok}[1]{\textcolor[rgb]{0.56,0.35,0.01}{{#1}}}
% \newcommand{\AlertTok}[1]{\textcolor[rgb]{0.94,0.16,0.16}{{#1}}}
% \newcommand{\FunctionTok}[1]{\textcolor[rgb]{0.00,0.00,0.00}{{#1}}}
% \newcommand{\RegionMarkerTok}[1]{{#1}}
% \newcommand{\ErrorTok}[1]{\textbf{{#1}}}
% \newcommand{\NormalTok}[1]{{#1}}

%espresso style colors
% \usepackage{framed}
% \definecolor{shadecolor}{RGB}{42,33,28}
% \newenvironment{Shaded}{\begin{snugshade}}{\end{snugshade}}
% \newcommand{\KeywordTok}[1]{\textcolor[rgb]{0.26,0.66,0.93}{\textbf{{#1}}}}
% \newcommand{\DataTypeTok}[1]{\textcolor[rgb]{0.74,0.68,0.62}{\underline{{#1}}}}
% \newcommand{\DecValTok}[1]{\textcolor[rgb]{0.27,0.67,0.26}{{#1}}}
% \newcommand{\BaseNTok}[1]{\textcolor[rgb]{0.27,0.67,0.26}{{#1}}}
% \newcommand{\FloatTok}[1]{\textcolor[rgb]{0.27,0.67,0.26}{{#1}}}
% \newcommand{\CharTok}[1]{\textcolor[rgb]{0.02,0.61,0.04}{{#1}}}
% \newcommand{\StringTok}[1]{\textcolor[rgb]{0.02,0.61,0.04}{{#1}}}
% \newcommand{\CommentTok}[1]{\textcolor[rgb]{0.00,0.40,1.00}{\textit{{#1}}}}
% \newcommand{\OtherTok}[1]{\textcolor[rgb]{0.74,0.68,0.62}{{#1}}}
% \newcommand{\AlertTok}[1]{\textcolor[rgb]{1.00,1.00,0.00}{{#1}}}
% \newcommand{\FunctionTok}[1]{\textcolor[rgb]{1.00,0.58,0.35}{\textbf{{#1}}}}
% \newcommand{\RegionMarkerTok}[1]{\textcolor[rgb]{0.74,0.68,0.62}{{#1}}}
% \newcommand{\ErrorTok}[1]{\textcolor[rgb]{0.74,0.68,0.62}{\textbf{{#1}}}}
% \newcommand{\NormalTok}[1]{\textcolor[rgb]{0.74,0.68,0.62}{{#1}}}

%kete style colors
% \newenvironment{Shaded}{}{}
% \newcommand{\KeywordTok}[1]{\textbf{{#1}}}
% \newcommand{\DataTypeTok}[1]{\textcolor[rgb]{0.50,0.00,0.00}{{#1}}}
% \newcommand{\DecValTok}[1]{\textcolor[rgb]{0.00,0.00,1.00}{{#1}}}
% \newcommand{\BaseNTok}[1]{\textcolor[rgb]{0.00,0.00,1.00}{{#1}}}
% \newcommand{\FloatTok}[1]{\textcolor[rgb]{0.50,0.00,0.50}{{#1}}}
% \newcommand{\CharTok}[1]{\textcolor[rgb]{1.00,0.00,1.00}{{#1}}}
% \newcommand{\StringTok}[1]{\textcolor[rgb]{0.87,0.00,0.00}{{#1}}}
% \newcommand{\CommentTok}[1]{\textcolor[rgb]{0.50,0.50,0.50}{\textit{{#1}}}}
% \newcommand{\OtherTok}[1]{{#1}}
% \newcommand{\AlertTok}[1]{\textcolor[rgb]{0.00,1.00,0.00}{\textbf{{#1}}}}
% \newcommand{\FunctionTok}[1]{\textcolor[rgb]{0.00,0.00,0.50}{{#1}}}
% \newcommand{\RegionMarkerTok}[1]{{#1}}
% \newcommand{\ErrorTok}[1]{\textcolor[rgb]{1.00,0.00,0.00}{\textbf{{#1}}}}
% \newcommand{\NormalTok}[1]{{#1}}
%end pandoc code hacks

% jodliterate colors
\usepackage{color}
\definecolor{shadecolor}{RGB}{248,248,248}
% j control structures 
\definecolor{keywcolor}{rgb}{0.13,0.29,0.53}
% j explicit arguments x y m n u v
\definecolor{datacolor}{rgb}{0.13,0.29,0.53}
% j numbers - all types see j.xml
\definecolor{decvcolor}{rgb}{0.00,0.00,0.81}
\definecolor{basencolor}{rgb}{0.00,0.00,0.81}
\definecolor{floatcolor}{rgb}{0.00,0.00,0.81}
% j local assignments
\definecolor{charcolor}{rgb}{0.31,0.60,0.02}
\definecolor{stringcolor}{rgb}{0.31,0.60,0.02}
\definecolor{commentcolor}{rgb}{0.56,0.35,0.01}
% primitive adverbs and conjunctions
%\definecolor{othercolor}{rgb}{0.56,0.35,0.01}   
\definecolor{othercolor}{RGB}{0,0,255}
% global assignments
\definecolor{alertcolor}{rgb}{0.94,0.16,0.16}
% primitive J verbs and noun names
\definecolor{funccolor}{rgb}{0.00,0.00,0.00}    

\usepackage{framed}
\newenvironment{Shaded}{}{}
\newcommand{\KeywordTok}[1]{\textcolor{keywcolor}{\textbf{{#1}}}}
\newcommand{\DataTypeTok}[1]{\textcolor{datacolor}{{#1}}}
%\newcommand{\DecValTok}[1]{\textcolor{decvcolor}{{#1}}}
\newcommand{\DecValTok}[1]{{#1}} 
\newcommand{\BaseNTok}[1]{\textcolor{basencolor}{{#1}}}
\newcommand{\FloatTok}[1]{\textcolor{floatcolor}{{#1}}}
\newcommand{\CharTok}[1]{\textcolor{charcolor}{\textbf{{#1}}}}
\newcommand{\StringTok}[1]{\textcolor{stringcolor}{{#1}}}
\newcommand{\CommentTok}[1]{\textcolor{commentcolor}{\textit{{#1}}}}
\newcommand{\OtherTok}[1]{\textcolor{othercolor}{{#1}}} 
\newcommand{\AlertTok}[1]{\textcolor{alertcolor}{\textbf{{#1}}}}
%\newcommand{\FunctionTok}[1]{\textcolor{funccolor}{{#1}}}
\newcommand{\FunctionTok}[1]{{#1}}
\newcommand{\RegionMarkerTok}[1]{{#1}}
\newcommand{\ErrorTok}[1]{\textbf{{#1}}}
\newcommand{\NormalTok}[1]{{#1}}

% headers and footers
\usepackage{fancyhdr}
\pagestyle{fancy}

\fancyhead{}
\fancyfoot{}

%\fancyhead[LE,RO]{\slshape \rightmark}
%\fancyhead[LO,RE]{\slshape \leftmark}
\fancyfoot[C]{\thepage}
%\headrulewidth 0.4pt
%\footrulewidth 0 pt

%\addtolength{\headheight}{\baselineskip}

%\lfoot{\emph{Analyze the Data not the Drivel}}
%\rfoot{\emph{\today}}

% subfigure handles figures that contain subfigures
%\usepackage{color,graphicx,subfigure,sidecap}
\usepackage{graphicx,sidecap}
\usepackage{subfigure}
\graphicspath{{./inclusions/}}

% floatflt provides for text wrapping around small figures and tables
\usepackage{floatflt}

% tweak caption formats 
\usepackage{caption} 
\usepackage{sidecap}
%\usepackage{subcaption} % not compatible with subfigure

\usepackage{rotating} % flip tables sideways

% complex footnotes
%\usepackage{bigfoot}

% weird logos \XeLaTeX
\usepackage{metalogo}

% source code listings
\usepackage{listings}

% long tables
% \usepackage{longtable}

\newcommand{\HRule}{\rule{\linewidth}{0.5mm}}

% map LaTeX cross references into PDF cross references
\usepackage[
            %dvips,
            colorlinks,
            linkcolor=blue,
            citecolor=blue,
            urlcolor=blue,   % magenta, cyan default        
            pdfauthor={John D. Baker},
            pdftitle={Analyze the Data not the Drivel},
            pdfsubject={Blog},
            pdfcreator={MikTeX+LaTeXe with hyperref package},
            pdfkeywords={blog,wordpress},
            ]{hyperref}
           
% custom colors
\definecolor{CodeBackGround}{cmyk}{0.0,0.0,0,0.05}    % light gray
\definecolor{CodeComment}{rgb}{0,0.50,0.00}           % dark green {0,0.45,0.08}
\definecolor{TableStripes}{gray}{0.9}                 % odd/even background in tables

\lstdefinelanguage{bat}
{morekeywords={echo,title,pushd,popd,setlocal,endlocal,off,if,not,exist,set,goto,pause},
sensitive=True,
morecomment=[l]{rem}
}

\lstdefinelanguage{jdoc}
{
morekeywords={},
otherkeywords={assert.,break.,continue.,for.,do.,if.,else.,elseif.,return.,select.,end.
,while.,whilst.,throw.,catch.,catchd.,catcht.,try.,case.,fcase.},
sensitive=True,
morecomment=[l]{NB.},
morestring=[b]',
morestring=[d]',
}

% latex size ordering - can never remember it
% \tiny
% \scriptsize
% \footnotesize
% \small
% \normalsize
% \large
% \Large
% \LARGE
% \huge
% \Huge
 
% listings package settings  
\lstset{%
  language=jdoc,                                % j document settings
  basicstyle=\ttfamily\footnotesize,            
  keywordstyle=\bfseries\color{keywcolor}\footnotesize,
  identifierstyle=\color{black},
  commentstyle=\slshape\color{CodeComment},     % colored slanted comments
  stringstyle=\color{red}\ttfamily,
  showstringspaces=false,                       
  %backgroundcolor=\color{CodeBackGround},       
  frame=single,                                
  framesep=1pt,                                 
  framerule=0.8pt,                             
  rulecolor=\color{CodeBackGround},   
  showspaces=false,
  %columns=fullflexible,
  %numbers=left,
  %numberstyle=\footnotesize,
  %numbersep=9pt,
  tabsize=2,
  showtabs=false,
  captionpos=b
  breaklines=true,                              
  breakindent=5pt                              
}

\lstdefinelanguage{JavaScript}{
  keywords={typeof, new, true, false, catch, function, return, null, catch, switch, var, if, in, while, do, else, case, break},
  ndkeywords={class, export, boolean, throw, implements, import, this},
  ndkeywordstyle=\color{darkgray}\bfseries,
  sensitive=false,
  comment=[l]{//},
  morecomment=[s]{/*}{*/},
  morestring=[b]',
  morestring=[b]"
}

% C# settings
\lstdefinestyle{sharpc}{
language=[Sharp]C,
basicstyle=\ttfamily\scriptsize, 
keywordstyle=\bfseries\color{keywcolor}\scriptsize,
framerule=0pt
}

% for source code listing longer than two use smaller font
\lstdefinestyle{smallersource}{
basicstyle=\ttfamily\scriptsize, 
keywordstyle=\bfseries\color{keywcolor}\scriptsize,
framerule=0pt
}

\lstdefinestyle{resetdefaults}{
language=jdoc,
basicstyle=\ttfamily\footnotesize,  
keywordstyle=\bfseries\color{keywcolor}\footnotesize,                                                               
framerule=0.8pt 
}

% APL UTF8 code points listed for lstlisting processing
\makeatletter
\lst@InputCatcodes
\def\lst@DefEC{%
 \lst@CCECUse \lst@ProcessLetter
  ^^80^^81^^82^^83^^84^^85^^86^^87^^88^^89^^8a^^8b^^8c^^8d^^8e^^8f%
  ^^90^^91^^92^^93^^94^^95^^96^^97^^98^^99^^9a^^9b^^9c^^9d^^9e^^9f%
  ^^a0^^a1^^a2^^a3^^a4^^a5^^a6^^a7^^a8^^a9^^aa^^ab^^ac^^ad^^ae^^af%
  ^^b0^^b1^^b2^^b3^^b4^^b5^^b6^^b7^^b8^^b9^^ba^^bb^^bc^^bd^^be^^bf%
  ^^c0^^c1^^c2^^c3^^c4^^c5^^c6^^c7^^c8^^c9^^ca^^cb^^cc^^cd^^ce^^cf%
  ^^d0^^d1^^d2^^d3^^d4^^d5^^d6^^d7^^d8^^d9^^da^^db^^dc^^dd^^de^^df%
  ^^e0^^e1^^e2^^e3^^e4^^e5^^e6^^e7^^e8^^e9^^ea^^eb^^ec^^ed^^ee^^ef%
  ^^f0^^f1^^f2^^f3^^f4^^f5^^f6^^f7^^f8^^f9^^fa^^fb^^fc^^fd^^fe^^ff%
  ^^^^20ac^^^^0153^^^^0152%
  ^^^^20a7^^^^2190^^^^2191^^^^2192^^^^2193^^^^2206^^^^2207^^^^220a%
  ^^^^2218^^^^2228^^^^2229^^^^222a^^^^2235^^^^223c^^^^2260^^^^2261%
  ^^^^2262^^^^2264^^^^2265^^^^2282^^^^2283^^^^2296^^^^22a2^^^^22a3%
  ^^^^22a4^^^^22a5^^^^22c4^^^^2308^^^^230a^^^^2336^^^^2337^^^^2339%
  ^^^^233b^^^^233d^^^^233f^^^^2340^^^^2342^^^^2347^^^^2348^^^^2349%
  ^^^^234b^^^^234e^^^^2350^^^^2352^^^^2355^^^^2357^^^^2359^^^^235d%
  ^^^^235e^^^^235f^^^^2361^^^^2362^^^^2363^^^^2364^^^^2365^^^^2368%
  ^^^^236a^^^^236b^^^^236c^^^^2371^^^^2372^^^^2373^^^^2374^^^^2375%
  ^^^^2377^^^^2378^^^^237a^^^^2395^^^^25af^^^^25ca^^^^25cb%  
  ^^00}
\lst@RestoreCatcodes
\makeatother

% custom lengths used within minipages
\newcommand{\minindent}{17pt}


\makeindex

\begin{document}

\subsection*{\href{https://bakerjd99.wordpress.com/2015/03/22/jod-update-version-0-9-97/}{JOD Update: Version 0.9.97*}}
\addcontentsline{toc}{subsection}{JOD Update: Version 0.9.97*}


\noindent\emph{Posted: 22 Mar 2015 22:53:00}
\vspace{6pt}

\captionsetup[floatingfigure]{labelformat=empty}
\begin{floatingfigure}[r]{0.10\textwidth}
\centering
\href{http://www.jsoftware.com/jwiki/Addons/general/jod}{\includegraphics[width=0.09\textwidth]{jodlittlecube.png}}
\label{fig:4946X0}
\end{floatingfigure} In the last year much has changed in the J world.

%\href{http://www.jsoftware.com/jwiki/Addons/general/jod}{\includegraphics{jodlittlecube.png}}

\begin{enumerate}
\itemsep1pt\parskip0pt\parsep0pt
\item
  There are new \href{http://www.jsoftware.com/stable.htm}{official J
  8.0x} builds for all supported platforms.
\item
  The \href{http://www.qt.io/download/}{QT} based IDE JDE has matured
  and is in widespread use.
\item
  The column oriented J database JD is drawing new users to J and
  enticing J veterans to reconsider how we use databases.
\item
  There is a small group of J system builders experimenting with
  additions, extensions and revisions of core J source code.
\end{enumerate}

In short, there are have been enough changes to revisit and update JOD.

JOD version
\texttt{0.9.97*}\footnote{The version number is \texttt{*}'ed because
you are always a point release from done!} is
the first JOD update in many years that mocks the god of software
compatibility. In particular:

\begin{enumerate}
\itemsep1pt\parskip0pt\parsep0pt
\item
  The syntax of the \texttt{jodhelp} verb has changed.
\item
  The \texttt{jodsource} addon no longer uses a zip file to distribute
  JOD dump files.
\item
  JOD online help will no longer be supported or updated.
\item
  Volume size is no longer checked before~creating new JOD dictionaries.
\item
  There is a new version of \texttt{jod.pdf.}
\end{enumerate}

\paragraph{\texorpdfstring{\texttt{jodhelp} changes (\#1, \#3 and
\#5):}{jodhelp changes (\#1, \#3 and \#5):}}

\texttt{jodhelp} has always been a \emph{kludge.} In programmer speak a
kludge is some half-baked facility added to a system after more
essential features have stabilized. The original versions of
\texttt{jodhelp} pointed at my rough notes. It was all the
``documentation'' I needed! Then others stared using JOD which resulted
in an ``evolved'' online version of my notes. I originally thought that
hosting my notes online would simultaneously serve user needs and cut
the amount of time I spent maintaining documentation. In retrospect this
\emph{wasn't even wrong!}

I used Google Documents to host my notes. If you've ever wondered why
completely free Google Documents hasn't obliterated expensive
\emph{Microsoft Word} or hoary old excellent \LaTeX\ I invite you
to maintain a set of \emph{long-duration-documents} with Google
Documents. During \texttt{jodhelp}'s online lifetime the basic internal
format of Google Documents changed in a \emph{screw-your-old-documents}
upgrade which forced me to spend days repairing broken hyper-links and
reformatting. I was not amused; you still get what you pay for!

I originally choose Google Documents because of its alleged global
accessibility. Sadly, Google Documents is now often blocked by corporate
and national firewalls. Even when it isn't blocked it renders like a dog
peeing on a fire hydrant. All these problems forced me to rewrite JOD
documentation with a completely reliable tool: good old-fashioned
\LaTeX. The result of my labors,
\href{https://app.box.com/shared/gajfu50gc0}{\texttt{jod.pdf}}, is now
distributed by the \texttt{joddocument} addon and is easily browsed with
\texttt{jodhelp}.

After \texttt{jod.pdf}'s appearance another irritant surfaced:
synchronizing \texttt{jod.pdf} and the online version. I tried using
\href{http://pandoc.org/}{pandoc} and
\href{http://daringfireball.net/projects/markdown/}{markdown} to
generate both the online and PDF versions from the same source files but
\texttt{jod.pdf} is too complex for \emph{not-to-fancy} portable
approaches. I was faced with a choice, lower my \texttt{jod.pdf}
standards, or get rid of something I never really liked. I opted to
drown a child and abandon online help. I don't expect a lot of mourners
at the funeral.

Using the new version of \texttt{jodhelp} requires installing the addon
\texttt{joddocument} and configuring a J PDF reader. It's also good idea
to define a JQT PF key to pop up JOD help with a keystroke. To configure
a J PDF reader edit the configuration file:

\begin{verbatim}
 ~config/base.cfg
\end{verbatim}

\noindent this file is directly available from the JQT
\texttt{Edit\textbackslash{}Configure} menu. \texttt{base.cfg} defines a
number of operating system dependent utilities. Make changes to the
systems you use, save your changes, and restart J. The following example
shows my \texttt{Win64} system settings.

\begin{lstlisting}[language=jdoc, frame=single,framerule=0pt,label=lst:scr4946X0]
 case. 'Win' do.   
   BoxForm=: 1   
   Browser=: 'c:/Program Files (x86)/Google/Chrome/Application/chrome.exe'   
   Browser_nox=: ''   
   EPSReader=: 'c:/program files/ghostgum/gsview/gsview64.exe'   
   PDFReader=: 'c:/uap/sumatra/SumatraPDF.exe'   
   XDiff=: 'c:/uap/WinMerge-2.14.0-exe/winmergeu.exe'   
   Editor=: 'c:/uap/notepad++/notepad++.exe %f'   
   Editor_nox=: '' 
\end{lstlisting}

I use
\href{http://www.sumatrapdfreader.org/free-pdf-reader.html}{SumatraPDF}
to read PDF files on Windows. It's a fast, lightweight, program that
efficiently renders \texttt{jod.pdf}. Good PDF readers are available for
all commonly used platforms.

To define JQT PK keys edit the configuration
file:\footnote{\texttt{userkeys.cfg} is only available for J 8.03 systems.
}

\begin{verbatim}
 ~config/userkeys.cfg 
\end{verbatim}

\noindent This file is also directly available from the
\texttt{Edit\textbackslash{}Configure} menu. My JOD specific PF keys
are:

\begin{lstlisting}[language=jdoc, frame=single,framerule=0pt,label=lst:scr4946X1]
 F3;1;Require JOD;require 'general/jod' 
 Shift+F3;1;JOD Help;jodhelp 0 
 F6;1;Dev Dicts;od cut 'joddev jod utils' [ 3 od '' 
 Shift+F6;1;Fit Dev Dicts;od cut 'jodfit joddev jod utils' [ 3 od '' 
 Ctrl+Shift+F6;1;Test Dev Dicts;od cut 'jodtest joddev jod utils' [ 3 od '' 
\end{lstlisting}

\noindent Pressing \texttt{Shift+F3} executes \texttt{jodhelp\ 0} which pops up
JOD help.

\paragraph{\texorpdfstring{\texttt{jodsource} changes
(\#2)}{jodsource changes (\#2)}:}

The \texttt{jodsource} addon is a collection of JOD dump scripts. Dump
scripts are serialized versions of binary JOD dictionaries. When
executed they merge objects into the current JOD put dictionary. I use
them primarily to move dictionaries around but they
\href{https://github.com/bakerjd99/joddumps}{have other uses} as well.
Prior to this version I distributed the three main JOD development dump
scripts, \texttt{joddev}, \texttt{jod}, and \texttt{utils} in one
compressed zip file to reduce the size of
\href{http://www.jsoftware.com/jwiki/JAL}{JAL} downloads.

The distributed script
\href{https://github.com/bakerjd99/jod/blob/master/jodsource/jodsourcesetup.ijs}{\texttt{jodsourcesetup.ijs}}
used the \texttt{zfiles} addon to extract these scripts and rebuild JOD
development dictionaries. This worked on 32 bit Windows systems but
failed elsewhere. J now runs on 32/64 bit Windows, Mac, Linux, IOS and
Android systems. To better support all these variants I eliminated the
\texttt{zfiles} dependency and pruned the JOD development dictionaries.
The result is a more portable and smaller \texttt{jodsource} addon.

\paragraph{Bye bye volume sizing (\#4):}

Early versions of JOD ran in the now bygone era of floppy disks. It was
possible to create \emph{many} JOD dictionaries on a single standard 800
kilobyte 3.5 inch floppy. Compared to modern \emph{porcine-ware} JOD,
which many J'ers consider a \emph{huge} system, is lithe and lean. In
floppy days it was important to check if there was enough space on a
floppy before creating another \emph{huge 48K} empty JOD dictionary.
This is a bit ridiculous today! If you don't have 48K free on whatever
device you are running you have far more serious problems than not being
able to create JOD dictionaries.

Volume sizing code remained in JOD for years until it started giving me
problems. Returning the size of very large network volumes can be
time-consuming and there are serious portability issues. Every operating
system calls different facilities to return volume sizes. Even worse,
security settings on corporate networks and cloud architectures
sometimes refuse to divulge national secrets like free byte counts.

To eliminate all these headaches this version of JOD no longer checks
volume size when the \texttt{FREESPACE} noun is zero. To restore the
previous behavior you have to edit the file

\begin{verbatim}
 ~addons/general/jod.ijs`
\end{verbatim}

\noindent and change the line \texttt{FREESPACE=:0} to whatever byte count you
want. Alternatively, you could
\emph{NGAF}\footnote{\emph{Not Give a Fuck!}
} and just assume you have 48K free on your terabyte size volumes.

\paragraph{Still to come:}

You may have surmised from JOD's version number that the system is still
not \emph{feature complete.} The JOD manual lists a few words that I am
planning to implement. I only develop JOD when I need something or I am
bored out of my mind at work and need a break. Such intermittent
motivators seldom insure project completion but I have found a new
reason to finish JOD. To list a book on
\href{http://www.goodreads.com/bakerjd99}{Goodreads} or Amazon you need
an \href{https://www.isbn-international.org/}{ISBN number}. The
\href{http://www.lulu.com/spotlight/bakerjd99}{hardcopy version} of the
JOD manual is a \emph{sort-of-published} book. To complete the
publishing process I need an ISBN. If I am going to bother with such
formalities I might as well complete the system the manual describes. So
there you have it a new software development motivator: vanity.

%\begin{center}\rule{0.5\linewidth}{\linethickness}\end{center}

%%begin{enumerate}
%\item
%  \hyperdef{}{fn3}{}
%
%  The version number is \texttt{*}'ed because you are always a point
%  release from done!\hyperref[fnref3]{↩}
%\item
%  \hyperdef{}{fn1}{}
%
%  \texttt{userkeys.cfg} is only available for J 8.03
%  systems.\hyperref[fnref1]{↩}
%\item
%  \hyperdef{}{fn2}{}
%
%  \emph{Not Give a F\%\&k!}\hyperref[fnref2]{↩}
%\end{enumerate}

%\captionsetup[floatingfigure]{labelformat=empty}
%\begin{figure}[htbp]
%\begin{floatingfigure}[l]{0.25\textwidth}
%\centering
%\includegraphics[width=0.23\textwidth]{jodlittlecube.png}
%\caption{~~~IMCAPTION~~~}
%\label{fig:4946X0}
%\end{floatingfigure}
%\end{figure}


%\end{document}